\chapter{Recuperación de Archivos Eliminados e Inodos Huérfanos}

Cuando un intruso advierte la posibilidad de ser detectado o concluye una fase de su ataque, suele ejecutar comandos de borrado (\texttt{rm -rf}) para suprimir sus herramientas y bitácoras intermedias. Sin embargo, en el sistema de archivos Ext2/Ext3, la eliminación de un archivo desasocia el inodo del directorio pero mantiene intactos los bloques de datos en el disco magnético hasta que son reutilizados por nuevas operaciones de escritura.

\begin{notaforense}[Mecánica Forense de Eliminación en Ext2: Inodos Huérfanos y Carving]
Cuando en un sistema Linux Ext2 se ejecuta la llamada al sistema \texttt{unlink()}:
\begin{enumerate}
    \item La entrada de directorio (\textit{dentry}) desasocia el nombre del archivo y el inodo.
    \item El contador de enlaces del inodo (\texttt{i_links_count}) disminuye a 0.
    \item Se registra la marca de tiempo de borrado (\texttt{i_dtime}) con el tiempo Epoch actual.
    \item Los mapas de bits de inodos y de bloques marcan esas estructuras como disponibles para el sistema operativo.
    \item \textbf{Los bloques de datos físicos NO se destruyen}: permanecen íntegros en los sectores del disco magnético (\textit{unallocated clusters}) hasta que otro archivo los sobrescribe.
\end{enumerate}
Mediante herramientas forenses como \texttt{ils} (para listar inodos desasociados con tiempo de borrado), \texttt{icat} (para volcar directamente los bloques asociados a un inodo huérfano) y técnicas de tallado de datos (\textit{file carving} por firmas de cabecera y pie), el perito puede resucitar archivos eliminados meses atrás si el bloque no fue reutilizado.
\end{notaforense}

\section{Recuperación de Inodos Huérfanos y Datos No Asignados (\textit{Carving})}


 Se recuperó el contenido de archivos en el disco duro que fueron
borrados por cualquier método y que no han sido sobreescritos aun. Se
usaron diversas técnicas de recuperación de datos entre ellas el barrido directo 
de todos los bits en el disco. Gran parte de la
información borrada no pudo extraerse ya que el servidor tenía un uso
muy intenso y el remplazo de los bloques de datos es casi inmediato. De los
archivos que no se pueden recuperar por medios estándar, se obtuvo el
nombre almacenado del archivo y se hizo un análisis de la cantidad informática llamada entropía 
de los bits en
los bloques del disco duro, con el fin de poder concatenar datos dispersos y
reconstruir algún archivo que pudiera ser de interés. Esta última
prueba fue realizada en ciertos archivos, pero su gran consumo de tiempo la descartó como
método de análisis general.


 Archivos sospechosos de entre 300,000 recuperados.

\vspace{1em}
\vspace{0.8em}\hrulefill\vspace{0.8em}
{\footnotesize
\begin{lstlisting}[style=terminal]
r / r	/dev/zqft4/algo	          2003.05.10 00:05:03 (cst)	2003.05.10 00:05:03 (cst)	
                                2003.05.10 00:05:09 (cst)	0	0	0	532496 (realloc)

s / s	   	/dev/log              2003.05.08 22:23:03 (cst)	   	2003.05.08 22:23:03 (cst)	   	
                                2003.05.10 03:00:57 (cst)	   	0	   	0	   	0
                                305889

d / r	   	/tmp/.sxp             2003.04.08 13:02:52 (cst)	   	2003.05.09 12:25:06 (cst)	   	
                                2003.04.23 22:34:05 (cst)	   	475738	   	0	   	0	   	
                                130441 (realloc)

r / r  	    	sh                 2003.04.15 21:39:50 (cst)  	    	2003.04.26 16:00:32 (cst)  	    	
                                2003.04.15 22:29:00 (cst)  	    	466184  	    	523  	    	
                                523  	   	786702 (realloc)

f / r  	    	mc.pipe.12891      2003.03.08 12:13:10 (cst)  	    	2003.05.08 03:01:44 (cst)  	    	
                                2003.04.11 20:55:39 (cst)  	    	69  	    	500  	    	100  	   	
                                130461 (realloc)
\end{lstlisting}
}
\vspace{0.8em}\hrulefill\vspace{0.8em}
\vspace{1em}

 Los archivos anteriores son logs del sistema y ejecutables de \textbf{IRC}. Lo
que es de destacar es la existencia de un archivo llamado \texttt{algo}.
Parece ser que \texttt{algo} fue una prueba de control. La palabra \texttt{algo}
está en español por lo que podemos decir que la instalación fue realizada
por un hispanoparlante. El archivo \texttt{/tmp/.sxp} fue recuperado e identificado
como parte del paquete de explotación contra el servicio \texttt{sendmail} en el directorio \texttt{/ex}.\\
 

 \texttt{mc.pipe.12891} es un archivo de tubería temporal creado por el administrador de archivos \textit{Midnight Commander} (\texttt{mc}), propiedad del usuario \texttt{duende}, y contiene referencias de un \textit{script} ejecutado desde la interfaz de dicho usuario.\\

 Su contenido es el siguiente:
 
\vspace{1em}
\vspace{0.8em}\hrulefill\vspace{0.8em}
{\footnotesize
\begin{lstlisting}[style=terminal]
linux
sysinit
dnsquery
identd
plogd
db
sh
.tree
screen
popclean
ecko
\end{lstlisting}
}
\vspace{0.8em}\hrulefill\vspace{0.8em}
\vspace{1em}

 Esto demuestra de manera incontrovertible que el 11 de abril, el usuario \texttt{duende}
ejecutó desde \texttt{Midnight Commander} varios de los programas maliciosos que
estamos estudiando. Recordemos que estos son archivos eliminados y
recuperados pericialmente para esta investigación.

\begin{retoforense}[Recuperación de Archivos Temporales de Aplicaciones]
El archivo \texttt{mc.pipe.12891} documentó la lista de procesos simulados del malware. ¿Por qué utilidades interactivas de consola como \texttt{mc}, \texttt{screen} o editores como \texttt{vim} dejan artefactos residuales en \texttt{/tmp} o \texttt{/var/tmp} y de qué manera un perito forense puede utilizarlos para reconstruir la interacción del usuario cuando el atacante ha eliminado su archivo \texttt{.bash_history}?
\end{retoforense}

\begin{perspectiva20}[Perspectiva 20 Años Después: File Carving en Discos SSD Modernos y TRIM]
En 2004, la recuperación de archivos borrados sobre discos magnéticos IDE/SCSI era una técnica pericial con altísima tasa de éxito, debido a que los bloques permanecían físicamente intactos hasta su sobreescritura accidental. En la actualidad, con el predominio absoluto de las unidades de estado sólido (SSD NVMe) y el comando \textit{TRIM / Deallocate}, cuando el sistema operativo borra un archivo, el controlador del SSD pone a cero inmediatamente los bloques lógicos para optimizar los ciclos de escritura. En dispositivos modernos con TRIM activo, el \textit{file carving} estático devuelve únicamente ceros (\texttt{0x00}), obligando a la disciplina forense contemporánea a priorizar la captura de memoria RAM volátil y copias de seguridad en la nube.
\end{perspectiva20}

